% Weekly update, 23 September 2026.
% Design follows docs/slides/weekly_update.tex: default Beamer, 4:3,
% Latin Modern, blue titles/bullets, white background and original navigation.
% Build: pdflatex -interaction=nonstopmode -halt-on-error weekly_update_20260923.tex
\documentclass{beamer}
\usepackage[utf8]{inputenc}
\usepackage{lmodern}
\usepackage{graphicx}
\usepackage{booktabs}
\usepackage{array}
\usepackage{tabularx}
\usepackage{hyperref}
\usepackage{listings}
\usepackage{xcolor}

\definecolor{codebg}{RGB}{247,247,247}
\definecolor{codeborder}{RGB}{220,220,220}
\lstdefinestyle{json}{basicstyle=\ttfamily\small,backgroundcolor=\color{codebg},
  frame=single,rulecolor=\color{codeborder},breaklines=true,showstringspaces=false}
\newcommand{\src}[1]{\par\vfill\begin{minipage}{0.67\textwidth}\tiny\color{gray}Source: #1.\end{minipage}\par}
\newcommand{\scope}[1]{{\small\textit{#1}}\par\vspace{0.5em}}
\newcolumntype{L}[1]{>{\raggedright\arraybackslash}p{#1}}
\renewcommand{\tabularxcolumn}[1]{>{\raggedright\arraybackslash}p{#1}}
\pdfstringdefDisableCommands{\def\translate#1{#1}}
\setlength{\tabcolsep}{5pt}

\title{LLM-based Power-System Agent\\Weekly Update}
\subtitle{Fault Corpora, DAgger R2\\and HIF-Conditioned Meter Repair}
\author{Project Team}
\date{September 23, 2026}

\begin{document}
\frame{\titlepage}

% 2
\begin{frame}{This Week's Changes}
  \begin{itemize}
    \setlength\itemsep{0.7em}
    \item \textbf{Corpora:} aligned sensor noise and covariance; introduced physical HIF resistance and consistent network references.
    \item \textbf{Completed DAgger round:} R2 achieved \textbf{141/160}, compared with R1's \textbf{140/160} on the same development roots.
    \item \textbf{Revised runtime and audit:} the unchanged R2 checkpoint reached \textbf{157/160 correct final states}, of which \textbf{153 terminated successfully}.
    \item \textbf{Current protocol:} SCADA-only discovery is implemented. The new physical corpora and this restricted protocol do not yet have a trained-model score.
  \end{itemize}
  \src{S1, S3--S5, S7. Completed scores use the September 21--22 auxiliary-assisted experiment}
\end{frame}

% 3
\begin{frame}{Noise, WLS and Audit Consistency}
  \begin{itemize}
    \item Draw noise with the covariance supplied to WLS:
      \[\sigma_V=0.001\ \mathrm{pu},\qquad \sigma_{P,Q}=0.01\ \mathrm{pu},\qquad
        \epsilon\sim\mathcal N(0,R).\]
    \item Detect an anomaly if either test triggers:
      \[J=e^\mathsf{T}R^{-1}e\geq\chi^2_{\nu,0.99}
        \quad\text{or}\quad \max_i |r_i^N|\geq4.\]
    \item A healthy control keeps its \textbf{original noisy acquisition}; the audit no longer requires removal of ordinary sensor noise.
    \item Mixed HIF/meter recovery is checked against the \textbf{HIF-present acquisition}, restoring only the gross-error targets.
  \end{itemize}
  \small Healthy non-target readings must remain unchanged. All processes use a \textbf{40-action} episode limit.
  \src{S2, S6, S8. Sensor-level sigmas; topology projections propagate covariance. Overall dual-alarm calibration is separate}
\end{frame}

% 4
\begin{frame}{Fault Injection Settings}
  \small
  \renewcommand{\arraystretch}{1.15}
  \begin{tabularx}{\textwidth}{L{2.55cm}X}
    \toprule
    Family & Main experimental setting \\
    \midrule
    Measurement & Gross offsets $\geq10\sigma$ in current DAgger. Multiple-meter cases have 2--5 targets. \\
    Parameter & One line, R/X/RX. Physical/reported factor in $[0.1,0.5]$ or $[2,5]$. \\
    Topology & Binary breaker-status error and circuit re-solve; require observable impact. \\
    HIF & Phase-to-ground resistor, 100--1000 $\Omega$ on 69-kV lines; 10 scans. \\
    Unbalance & Dirichlet$(3,3,3)$ phase allocation, preserving total P/Q; achieved severity is measured by VUF. \\
    Harmonic & Source-bus voltage THD 10--20\%; orders 5, 7, 11, 13, 17 and 19. \\
    \bottomrule
  \end{tabularx}
  \par\vspace{0.6em}
  \small Mixed meter+parameter/topology/HIF cases add one power-channel offset of signed 0.10--0.30 pu.
  \src{S2, S8. Main training settings; lower-severity sensitivity populations are separate}
\end{frame}

% 5
\begin{frame}{HIF Resistance on Physical Voltage Bases}
  \[Z_{\rm base}=\frac{V_{\rm LL,base}^2}{S_{\rm 3\phi,base}},\qquad
    R_{f,\rm pu}=\frac{R_{f,\Omega}}{Z_{\rm base}},\qquad S_{\rm base}=100\ \mathrm{MVA}.\]
  \begin{center}
    \begin{tabular}{lrrr}
      \toprule
      Model voltage & $Z_{\rm base}$ ($\Omega$) & 100 $\Omega$ (pu) & 500 $\Omega$ (pu) \\
      \midrule
      69 kV & 47.61 & 2.10 & 10.50 \\
      13.8 kV & 1.9044 & 52.51 & 262.55 \\
      18 kV & 3.24 & 30.86 & 154.32 \\
      138 kV & 190.44 & 0.53 & 2.63 \\
      \bottomrule
    \end{tabular}
  \end{center}
  \begin{itemize}
    \item IEEE-14 realization: buses 1--5 at 69 kV, bus 8 at 18 kV, remaining buses at 13.8 kV.
    \item Main HIF corpus: seven 69-kV lines; phases A/B/C; fault position 0.25--0.75 of line length.
    \item IEEE-57 uses a declared 138/69-kV reconstruction. The DAgger results in this deck are \textbf{IEEE-14}.
  \end{itemize}
  \src{S1, S10. Main DAgger DSS uses a 1-kV equivalent with explicit local-pu/ohm conversion; no arcing model}
\end{frame}

% 6
\begin{frame}{Physical Corrections to the Corpora}
  \begin{itemize}
    \setlength\itemsep{0.55em}
    \item \textbf{Network convention:} align OpenDSS exports with the balanced WLS \texttt{ybus} shunt convention and local voltage/current bases.
    \item \textbf{Reference state:} replace the OPF reference used as physical truth with a paired OpenDSS solve at the same operating point and dispatch.
    \item \textbf{September 23 generator fix:} restore declared reactive limits after generator kW writes. Previously, synchronous condensers lost reactive capability and created a fault-family cue.
    \item Regenerate six HIF corpora with the same recipes and seeds. The completed audited pool increased from \textbf{118 to 131} detectable windows.
  \end{itemize}
  \small The completed ``20260923b'' rebuild also aligns both generators at a $10^{-8}$ solve tolerance. Earlier DAgger scores retain their original physical model.
  \src{S1. New HIF/unbalance corpora and detection audits completed; five legacy NLM top-1 misses remain}
\end{frame}

% 7
\begin{frame}{WLS Admission and Detection Limits}
  \small
  Admission margin: $J>1.25\tau_{\chi^2}$ or $\max_i|r_i^N|\geq5$.
  \begin{center}
    \renewcommand{\arraystretch}{1.15}
    \begin{tabular}{lrrr}
      \toprule
      HIF source & Windows & Prior admitted & New admitted \\
      \midrule
      Main training & 84 & 25 & 27 \\
      Main validation & 21 & 7 & 8 \\
      Additional training & 252 & 69 & 77 \\
      Additional validation & 63 & 17 & 19 \\
      \midrule
      Total & 420 & 118 & \textbf{131} \\
      \bottomrule
    \end{tabular}
  \end{center}
  \begin{itemize}
    \item New admitted HIF bands: \textbf{113} at 100--200 $\Omega$, \textbf{18} at 200--500 $\Omega$, \textbf{0} at 500--1000 $\Omega$.
    \item Unbalance: 160/440 admitted; 0/60 healthy-control alarms at the admission margin (September 23b audit).
    \item Weak HIFs remain in evaluation. Three of the 131 admitted HIF windows also alarm in paired healthy data.
  \end{itemize}
  \src{S1. New HIF counts use September 23 physics; completed DAgger used the prior 118-window pool}
\end{frame}

% 8
\begin{frame}{Completed DAgger Training Run}
  \scope{Gemma 4 12B; September 21--22 auxiliary-assisted run}
  \begin{itemize}
    \setlength\itemsep{0.55em}
    \item Initial expert aggregate: \textbf{5,021 rows from 506 roots}; BC0 trained on \textbf{3,784 rows}.
    \item R2 collection: \textbf{122 episodes}, 1,099 visited states and \textbf{660 eligible learner-state labels}.
    \item R2 mixture: \textbf{660 replay + 660 new rows}; \textbf{330 training steps} completed.
    \item Evaluation: the same \textbf{160 development roots}, fixed seed 20260912 and a \textbf{40-action} limit.
    \item HIF fit: $7\times9$ search, up to 10 scans. Policy generation: 32,768 input / 256 output tokens.
  \end{itemize}
  \small HIF was \textbf{flagged} in this run. Filtering its inputs for WLS detectability did not remove the initial HIF hint or auxiliary telemetry.
  \src{S3, S8. Rows, visited states and physical roots are different units}
\end{frame}

% 9
\begin{frame}{Original DAgger Results}
  \scope{Same 160 roots and original runtime/audit}
  \begin{center}
    \renewcommand{\arraystretch}{1.25}
    \begin{tabular}{lrr}
      \toprule
      Policy & Truth-audited success & Rate \\
      \midrule
      BC0 & 132/160 & 82.50\% \\
      R1 & 140/160 & 87.50\% \\
      R2 & \textbf{141/160} & \textbf{88.125\%} \\
      Expert & 142/160 & 88.75\% \\
      \bottomrule
    \end{tabular}
  \end{center}
  \begin{itemize}
    \item R1 $\rightarrow$ R2: two gains and one regression; net \textbf{+1 root / +0.625 percentage points}.
    \item Gains: one multi-meter case and one parameter case. Regression: one mixed meter/parameter case.
    \item R2's 141 successes comprise 56 resolved outcomes and 85 truth-correct handoffs. Waveform diagnosis does not remove the physical disturbance.
  \end{itemize}
  \src{S3. Development evaluation, not a new independent test set}
\end{frame}

% 10
\begin{frame}{Reevaluation with the Revised Code}
  \scope{Unchanged R2 weights, roots, noise, thresholds and action budget}
  \small
  \begin{center}
    \renewcommand{\arraystretch}{1.3}
    \begin{tabular}{lrrr}
      \toprule
      & Original & Revised & Change \\
      \midrule
      Expert: audit-correct & 142/160 & \textbf{158/160} & +16 \\
      R2: audit-correct & 141/160 & \textbf{157/160} & +16 \\
      R2: correct \emph{and terminal} & 141/160 & \textbf{153/160} & +12 \\
      \bottomrule
    \end{tabular}
  \end{center}
  \begin{itemize}
    \setlength\itemsep{0.5em}
    \item \textbf{Eight healthy controls:} corrected audit reference; their noisy measurements were already preserved.
    \item \textbf{Eight mixed HIF cases:} actual meter repairs became possible after conditioning on the HIF estimate.
    \item Four repaired R2 cases did not complete a valid handoff. All 158 expert successes terminated.
  \end{itemize}
  \vspace{0.3em}
  \textbf{The +16 reflects runtime/audit changes, not another round of model training.}
  \src{S4, S5. Auxiliary-assisted results on the frozen September 21 inputs}
\end{frame}

% 11
\begin{frame}{Masking the HIF Contribution}
  \scope{Implemented and evaluated with auxiliary phase/current evidence}
  Estimate the HIF and simulate the same operating point with and without it:
  \[\widehat{\Delta z}_{\rm HIF}
    =\widehat z_{\rm HIF\ present}-\widehat z_{\rm HIF\ absent}.\]
  Then run balanced WLS on a temporary conditioned vector:
  \[\boxed{z_{\rm conditioned}=z_{\rm observed}-\widehat{\Delta z}_{\rm HIF}}.\]
  \begin{itemize}
    \item Retain every channel and the original covariance $R$.
    \item Preserve raw observations; do not zero residuals or delete affected meters.
    \item A coexisting meter error remains in the conditioned vector, including on a channel affected by HIF.
  \end{itemize}
  \small The current SCADA scan is excluded from HIF fitting. $R$ does not include HIF-fit uncertainty; conditional chi-square calibration remains unestablished.
  \src{S6. The physical HIF remains on the network}
\end{frame}

% 12
\begin{frame}{Continuing into Meter Repair}
  \begin{enumerate}
    \setlength\itemsep{0.35em}
    \item Localize and fit HIF from the auxiliary measurements and other scans.
    \item Run conditioned WLS on the current acquisition.
    \item Obtain meter context; require a sparse discrepancy exceeding \textbf{5 measurement sigmas} outside the sampled prediction range.
    \item Replace supported targets with their \textbf{HIF-present prediction}; preserve every other reading.
    \item Verify the candidate, commit only a supported repair, then hand off for the remaining physical fault.
  \end{enumerate}
  \vspace{0.6em}
  \small
  \textbf{Guards:} prediction-range width $\leq2\sigma$; fresh acquisition/model binding; no thermal or topology violation. A changed voltage-limit violation cannot use the partial-repair exception.
  \src{S6. The sampled sensitivity range is not a statistical confidence interval}
\end{frame}

% 13
\begin{frame}{HIF-Conditioned Repair Verification}
  \scope{September 22 expert replays on eight frozen physical roots}
  \small
  \begin{center}
    \renewcommand{\arraystretch}{1.25}
    \begin{tabularx}{\textwidth}{Xr}
      \toprule
      Test & Result \\
      \midrule
      Original HIF + meter: repair and handoff & \textbf{8/8} \\
      Added $10\sigma$ error on strongest HIF channel: detected & 8/8 \\
      Added-overlap cases: committed and audit-correct & \textbf{6/8} \\
      Paired HIF-only controls: diagnosis, no meter writes & 8/8 \\
      Non-target writes across the 24 episodes & \textbf{0} \\
      \bottomrule
    \end{tabularx}
  \end{center}
  On the original mixed roots, \textbf{conditioned} statistics changed as follows:
  \[J:\ 246.1\text{--}847.7\ \longrightarrow\ 77.2\text{--}114.1,\qquad \tau_J=130.0\]
  \[\max|r^N|:\ 12.24\text{--}27.46\ \longrightarrow\ 2.24\text{--}2.94.\]
  \small Two overlap stress cases changed voltage readings that still violated limits; both were rolled back and retained as failures.
  \src{S6. Paired variants share eight roots; these are not 24 independent physical events}
\end{frame}

\begin{frame}{Example: HIF and Meter Error Overlap}
  \scope{September 22 auxiliary-assisted replay; frozen pre-fix physics}
  \small At bus 1's active-power injection meter ($\sigma=0.01$ pu):
  \begin{center}
    \renewcommand{\arraystretch}{1.2}
    \begin{tabular}{lr}
      \toprule
      Quantity & pu \\
      \midrule
      Predicted HIF-absent value & 1.948320 \\
      Predicted HIF-present value & 2.100483 \\
      Estimated HIF contribution & $+0.152163$ \\
      Added meter error & $-0.100000$ \\
      Observed corrupted value & 1.998634 \\
      Conditioned value & 1.846471 \\
      \bottomrule
    \end{tabular}
  \end{center}
  \[1.998634-0.152163=1.846471\ \mathrm{pu}.\]
  The meter bias survives compensation. Repair restores the raw target to \textbf{2.100483 pu}, preserving the HIF contribution.
  \[J_{\rm conditioned}:\ 639.45\ \longrightarrow\ 78.47.\]
  \src{S6, controlled-overlap root 5fb25759717c. Both faulty meters repaired; no other readings changed}
\end{frame}

% 14
\begin{frame}{Remaining R2 Failures}
  \scope{Revised auxiliary-assisted evaluation; unchanged checkpoint}
  \begin{itemize}
    \setlength\itemsep{0.6em}
    \item \textbf{Four correct but nonterminal states:} HIF meter repairs committed, followed by unsupported handoff request strings. Repetition stopped the episodes after 16--18 actions.
    \item \textbf{Three wrong final states:} one remaining multi-meter error, one wrong parameter branch, and one mixed case with wrong-branch and healthy-meter edits.
  \end{itemize}
  \small
  \begin{center}
    \renewcommand{\arraystretch}{1.15}
    \begin{tabular}{lrr}
      \toprule
      Behavioral counter & Original R2 & Revised R2 \\
      \midrule
      Invalid actions & 21 & 65 \\
      Episodes with a loop & 2 & 6 \\
      Rejected commit attempts & 15 & 15 \\
      False-finalization counter & 16 & 26 \\
      \bottomrule
    \end{tabular}
  \end{center}
  Original 16 finalizations completed but failed the audit. Revised 26 were rejected without mutation. Accepted wrong edits are audited separately.
  \src{S3, S4. Repair accuracy improved while protocol reliability remained a limitation}
\end{frame}

% 15
\begin{frame}{Current SCADA-Only Discovery Protocol}
  \begin{itemize}
    \setlength\itemsep{0.65em}
    \item Every episode starts with \textbf{balanced WLS}; no seeded HIF, unbalance or harmonic flag.
    \item Agent inputs: network model, SCADA readings, covariance and permitted SCADA history. Auxiliary phase/current and harmonic tools are blocked.
    \item The existing HIF fit and compensation depend on auxiliary evidence and are \textbf{disabled} in this profile.
    \item The eight frozen mixed roots still triggered both WLS alarms in the initial check, with identical actions after hidden labels and auxiliary data were changed.
  \end{itemize}
  \vspace{0.5em}
  \textbf{Full SCADA-only HIF identification and repair performance is not yet measured.} Earlier scores cannot be transferred to this evidence setting.
  \src{S7. Initial-alarm check on September 22 frozen inputs; no new training or full evaluation}
\end{frame}

% 16
\begin{frame}{Next Experiments}
  \begin{enumerate}
    \setlength\itemsep{0.9em}
    \item \textbf{Freeze the rebuilt experiment.} Bind the completed HIF/unbalance corpora to their source, model and detector settings before collecting new trajectories.
    \item \textbf{Collect SCADA-only expert traces.} Keep weak and unresolved cases in evaluation; exclude unsupported teacher actions from training.
    \item \textbf{Measure the new baseline.} Evaluate expert and frozen R2 on matched inputs before attributing gains to retraining.
    \item \textbf{Develop SCADA-only HIF conditioning.} Test identifiable HIF hypotheses against competing meter errors, then validate overlap and model-uncertainty limits.
  \end{enumerate}
  \src{S1, S7. Planned work; the completed results in this deck use earlier frozen configurations}
\end{frame}

\appendix
% 17
\begin{frame}{Appendix: Success by Fault Family}
  \footnotesize
  \renewcommand{\arraystretch}{1.08}
  \begin{tabularx}{\textwidth}{Xrrrrr}
    \toprule
    Family & $N$ & R1 & R2 & Revised R2 & Rev. expert \\
    \midrule
    Measurement & 16 & 16 & 16 & 16 & 16 \\
    Multiple measurements & 12 & 10 & 11 & 11 & 12 \\
    Parameter & 16 & 14 & 15 & 15 & 14 \\
    Measurement + parameter & 16 & 16 & 15 & 15 & 16 \\
    Topology & 16 & 16 & 16 & 16 & 16 \\
    Measurement + topology & 12 & 12 & 12 & 12 & 12 \\
    HIF & 16 & 16 & 16 & 16 & 16 \\
    Measurement + HIF & 8 & 0 & 0 & 8 & 8 \\
    Unbalance & 16 & 16 & 16 & 16 & 16 \\
    Harmonic & 16 & 16 & 16 & 16 & 16 \\
    No-error controls & 8 & 8 & 8 & 8 & 8 \\
    Healthy telemetry controls & 8 & 0 & 0 & 8 & 8 \\
    \midrule
    Total & 160 & 140 & 141 & 157 & 158 \\
    \bottomrule
  \end{tabularx}
  \par\vspace{0.5em}
  \small Revised R2's mixed-HIF row contains four nonterminal states. Its correct-and-terminal total is \textbf{153/160}.
  \src{S3--S5. Final-state truth audit on the same 160 roots; auxiliary-assisted configurations}
\end{frame}

% 18
\begin{frame}{Appendix: Interpretation and Sources}
  \small
  \begin{itemize}
    \setlength\itemsep{0.55em}
    \item \textbf{Correct final state:} passes the task-specific truth audit; may be nonterminal. \textbf{Correct and terminal:} additionally completes a valid resolution or handoff.
    \item \textbf{Audited completion:} a stricter stored certificate metric; revised R2 scores 136/160. It is different from 153/160 correct terminal tasks.
    \item \textbf{Physical removal:} none of these HIF handoffs clears the physical HIF. Diagnostic success and meter repair are separate outcomes.
    \item \textbf{Scope:} synthetic resistive HIFs, selected detectable roots, repeated development evaluation, and a new physical-model revision awaiting a matched evaluation.
  \end{itemize}
  \vspace{0.5em}
  \footnotesize
  S1: physical-corpus revision and audits. S2: injection generators.\\
  S3: original R2 report. S4/S5: revised R2/expert reports.\\
  S6: HIF continuation method and 24-episode replay.\\
  S7: SCADA-only protocol. S8: pipeline/training receipts.\\
  S10: IEEE-14/57 voltage-base conventions.\\[0.4em]
  Exact paths, hashes, slide mappings and presenter notes are included with the LaTeX source.
\end{frame}
\end{document}
